\subsection{帰納法の準備}

$R$は可換環とする．

\begin{definition}
  \label{def:CartanMatrix.rev}
  \lean{CartanMatrix.rev}
  \leanok
  $X \in \M_n(R)$に対し，$\rev X$で$X$の$(i, j)$成分と$(n - i + 1, n - j + 1)$成分を入れ替えた行列を表す．
\end{definition}

\begin{definition}
  \label{def:CartanMatrix.ind_matrix}
  \lean{CartanMatrix.ind_matrix}
  \leanok
  $Y \in \M_n(R),\ a, b \in R$に対し，$\ind(Y, a, b) \in \M_{n+1}(R)$を以下のように定める：
  \begin{equation}
    \ind(Y, a, b)
    = \left(
      \begin{array}{cccc|c}
        & & & & 0 \\
        & Y & & & \vdots \\
        & & & & 0 \\
        & & & & b \\ \hline
        0 & \cdots & 0 & a & 2
      \end{array}
    \right).
  \end{equation}
\end{definition}

\begin{lemma}
  \label{lem:CartanMatrix.ind_det}
  \lean{CartanMatrix.ind_det}
  \leanok
  $X \in \M_{n+2}(R),\ Y \in \M_{n+1}(R),\ Z \in \M_n(R),\ a, b \in R$とし，$X = \ind(Y, a, b),\ Y_{n-1} = Z$とする．
  すなわち，以下が成り立っているとする：
  \begin{equation}
    X
    =\left(
      \begin{array}{cccc|c}
        & & & & 0 \\
        & Y & & & \vdots \\
        & & & & 0 \\
        & & & & b \\ \hline
        0 & \cdots & 0 & a & 2
      \end{array}
    \right)
    = \left(
      \begin{array}{ccc|c|c}
        & & & & 0 \\
        & Z & & * & \vdots \\
        & & & & 0 \\ \hline
        & * & & * & b \\ \hline
        0 & \cdots & 0 & a & 2
      \end{array}
    \right).
  \end{equation}
  このとき，以下の関係式が成り立つ：
  \begin{equation}
    \det X = 2 \det Y - ab \det Z.
  \end{equation}
\end{lemma}

\begin{proof}
  余因子展開すれば得られる．
\end{proof}

\begin{lemma}
  \label{lem:CartanMatrix.A_leadingSubmatrix}
  \lean{CartanMatrix.A_leadingSubmatrix}
  \leanok
  各Cartan行列$X_n$に対し，$k < n$次主小行列$(X_n)_k$は表~\ref{table:CartanMatrix.A_leadingSubmatrix}のようになる：
  \begin{table}[htbp]
    \centering
    \caption{Cartan行列の主小行列}
    \label{table:CartanMatrix.A_leadingSubmatrix}
    \begin{tabular}{c|ccccccccc}
      $X_n$ & $A_n$ & $B_n$ & $C_n$ & $D_n$ & $\rev D_n$ & $E_{n \le 8}$ & $F_4$ \\ \hline
      $(X_n)_k$ & $A_k$ & $A_k$ & $A_k$ & $A_k$ & $\rev D_k$ & $E_k$ & $B_3, A_{k \ne 3}$
    \end{tabular}
  \end{table}
\end{lemma}

\begin{proof}
  Lean内では各$(i, j)$成分に関して，定義を展開して確かめている．
  人が確かめるには，$X_n$の$k$次主小行列は，$X_n$のCoxeter-Dynkin図形から頂点$k+1, \ldots, n$を取り除いたグラフであるから，図~\ref{fig:positive_definite_graph}を見ればよい．
\end{proof}